\section{Later Scholastic Disputations and the Byzantine Tradition}

How can two pure spirits differ from each other if neither has a body? Picture many seals pressed into wax, St.~Bonaventure suggests in his \emph{Commentary on the Sentences} (\emph{In II Sent.}, d.~3, p.~1, a.~2, q.~3): ``neither can the seals be multiplied without the wax, nor is the wax numbered except because diverse seals are stamped upon it'' (\emph{nec sigilla plurificari possunt sine cera, nec cera numeratur, nisi quia fiunt in ea diversa sigilla}). Around questions like this---whether angels have spiritual matter, whether several angels share one species, how a pure spirit occupies place, and how finite minds behold the infinite God---the great Catholic schools after St.~Thomas Aquinas worked out distinct metaphysical architectures inside the same common faith, while the Byzantine East carried forward an unbroken patristic theology anchored in cosmic liturgy and deifying light.

\subsection{St.~Bonaventure and the Franciscan School}

St.~Bonaventure (1221--1274), writing his \emph{Commentary on the Second Book of Sentences} and his concise \emph{Breviloquium} roughly a decade before Aquinas composed the \emph{Prima Pars}, inherits Peter Lombard's synthesis (\emph{II Sent.}, dd.~2--11) and sets the classic Franciscan course on four contested points.

First, whereas Aquinas holds every angel to be a pure subsistent form free from all matter (\ST{I, q.~50, a.~2}), Bonaventure defends universal hylomorphism---a real composition of spiritual matter and form in every created spirit. Bonaventure rejects any corporeal, quantitative, or heterogeneous parts in angels; in \emph{Breviloquium}~II.6 he states plainly that an angel is \emph{substantia spiritualis et incorporea, ac per hoc simplex, omni carens dimensione quantitativa} (``a spiritual and incorporeal substance, and thereby simple, lacking all quantitative dimension''). Yet extending the word \emph{materia} broadly to ``every constitutive principle of potentiality'' (\emph{omne potentiale constitutivum}, \emph{In II Sent.}, d.~3, p.~1, a.~1, q.~1), he argues that an angel's real mutability across changing acts, its receptivity to illumination, its creaturely limitation, and its essential composition demand two distinct principles standing to one another as act and potency, hence as form and matter: \emph{ipsa substantia Angeli composita est ex materia et forma} (``the very substance of an angel is composed of matter and form''). Considered in itself prior to determination, matter is ``neither spiritual nor corporeal'' (\emph{nec est spiritualis, nec corporalis}, d.~3, p.~1, a.~1, q.~2, ad~3); united to spiritual form, it yields an incorporeal creature without collapsing the gulf between creature and the wholly immaterial God (citing Boethius, \emph{De Trinitate}~c.~2).

Second, spiritual matter opens a direct path to rejecting the Thomistic thesis that every angel constitutes a separate species (\ST{I, q.~50, a.~4}). In \emph{In II Sent.}, d.~3, p.~1, a.~2, q.~1, Bonaventure judges the claim \emph{tot species, quot individua} (``as many species as individuals'') to be an unwarranted \emph{praesumtio} (presumption) lacking cogent proof in Scripture or the Fathers, pointing out that Scripture assigns many angels to a single office and common work. He upholds instead as \emph{positio sobria et catholica} (``a sober and Catholic position'') that angels---either some or all---share a common specific nature and differ solely in personal number, because ``the love of charity exults in a multitude of good fellowship'' (\emph{amor caritatis exsultat in multitudine bonae societatis}). Just as human beings wear diverse faces and diverse mental traits within one species, so angels differ by their individualizing principles and innate properties, arising from the actual conjunction of spiritual matter and form like multiple seals stamped upon wax, while personal dignity follows principally from form (\emph{In II Sent.}, d.~3, p.~1, a.~2, q.~3).

Third, Bonaventure diverges from Aquinas on the state of the angels at creation and the interval of their trial. Taking over Peter Lombard's four primal attributes (\emph{II Sent.} d.~3)---simplicity of essence, personal distinction, natural image of the Trinity in memory, intellect, and will, and freedom of choice---Bonaventure pairs them in \emph{Breviloquium}~II.6 with four concomitant endowments: vigor in action (\emph{virtuositas}), readiness in ministry (\emph{officiositas}), keenness in discernment (\emph{perspicacitas}), and immutability after choice (\emph{immutabilitas post electionem}). Unlike Aquinas, who holds that angels were created in sanctifying grace (\ST{I, q.~62, a.~3}), Bonaventure teaches in \emph{Breviloquium}~II.7 that God created all angels naturally good, poised \emph{medios tamen inter se, summum bonum, et commutabile bonum, quod est creatura} (``midway between himself, the supreme good, and changeable good, which is the creature''): turning upward toward God raised them into grace and glory, whereas turning downward toward private excellence plunged Lucifer and his adherents into sin and immediate expulsion from the empyrean into dark air or hell. Presuming on his natural splendor and coveting private pre-eminence, Lucifer made himself his own origin and end; by natural immutability after choice he became instantly obstinate, his keen mind bent to deception and his vigor curtailed to sudden physical transmutations. Meanwhile the faithful spirits were instantly confirmed in glory, enjoying morning and evening knowledge and running within God wherever they are dispatched (\emph{quocumque mittantur, intra Deum currunt}, \emph{Breviloquium}~II.8), ordered across three triads answering to discernment (Thrones, Cherubim, Seraphim), vigor (Dominations, Virtues, Powers), and ministry (Principalities, Archangels, Angels---the last guarding ``lest those standing fall, and helping the fallen to rise'').

Three years after both Aquinas and Bonaventure died in 1274, this scholastic debate collided with ecclesiastical censure at Paris. On 7 March 1277, Bishop Stephen Tempier condemned 219 propositions taught in the Paris Arts faculty, several of which touched Thomistic angelology (documented alongside classic excerpts from Scotus and Cajetan in the 1885 Quaracchi \emph{Scholia} to Bonaventure's \emph{Sentences}). As Durandus of Saint-Pour\c{c}ain records, three articles struck at the species-uniqueness of pure forms: Article~81 censured the claim that because intelligences lack matter, God could not create several in one species; Article~96 censured the claim that God cannot multiply individuals under one species without matter; and Article~191 censured the claim that forms receive division only through matter. A fourth censure, Article~25, struck at the operational theory of angelic location: ``Separate substances are somewhere through operation---an error, if understood to mean that without operation a substance is not in place.'' Though Bishop Tempier's act bound only the diocese of Paris and his successor later revoked the articles touching Aquinas following Aquinas's canonization in 1323, those 1277 censures permanently shaped the next three centuries of scholastic angelology.

\subsection{Bl.~John Duns Scotus, Cajetan, and Francisco Su\'arez}

Working in the shadow of 1277, Bl.~John Duns Scotus (c.~1266--1308) forged a third metaphysical path distinct from both Aquinas and Bonaventure. Rejecting Bonaventure's spiritual matter in angels, Scotus still safeguarded God's freedom to multiply angels within one species through his celebrated doctrine of individuation (\emph{In II Sent.}, d.~3, q.~6): every creature is individuated by a positive individual entity superadded to the specific essence and formally distinct from it on the side of reality (\emph{ex parte rei formaliter distincta})---the individual difference commonly called \emph{haecceitas} (``thisness''), which contracts a common nature to a singular individual just as a specific difference contracts a genus. Applied to angels, \emph{haecceitas} allows immaterial spirits without any matter to belong to the same species while remaining genuinely distinct persons.

Scotus applied the same critical rigor to angelic location (\emph{In II Sent.}, d.~2, q.~6, nn.~11 and~15), rejecting Aquinas's view that an angel occupies place solely by exercising power upon a body (\ST{I, q.~52, a.~1}) and appealing directly to Paris 1277 Article~25. An angel is not necessarily in place by nature---God could create an angel without creating any corporeal universe at all, or outside every body---yet an angel possesses a passive aptitude to be definitively present in a containing corporeal place, first placed there by God and thereafter able to move itself from place to place without needing to act physically upon that spot. Two further hallmarks distinguish Scotus and his school across later theology: on angelic speech (\ST{I, q.~107, a.~1}), Scotus accounts for communication through an intelligible species actively impressed by the speaking angel upon the listener's intellect rather than bare volitional unveiling; and on Lucifer's first sin (\emph{Ordinatio}~II, d.~6, q.~2), Scotus himself identifies an immoderate love of his own beatitude (\emph{affectio commodi}, a spiritual self-indulgence or \emph{luxuria spiritualis}) unmoderated by justice (\emph{affectio iustitiae}), while a later Scotist line---building on Scotus's doctrine of the absolute primacy of Christ---came to link Lucifer's revolt with proud refusal to adore the revealed Incarnation.

Two centuries later, the great Dominican commentator Cardinal Thomas de Vio Cajetan (1469--1534) rallied to defend Aquinas's core theses---both \emph{tot species, quot individua} (\ST{I, q.~50, a.~4}) and presence in place by virtual contact (\ST{I, q.~52, a.~1})---against Scotist and nominalist objections. Yet Cajetan paired razor-sharp logic with disarming modesty before angelic mystery. Commenting on \ST{I, q.~52, a.~1}, Cajetan observes that all Catholic doctors agree an angel is in place definitively through virtual quantity and virtual contact rather than bodily mass, and that this contact constitutes a real relation of presence to place; yet he candidly adds that the exact foundation of that presence ``surpasses the capacity of human wit in its evidence, because we learn from the sacred Scriptures that angels are somewhere, yet it does not appear how they touch that place'' (\emph{Excedit enim evidentia eius humani ingenii capacitatem, propterea quod ex sacris eloquiis habemus, Angelos alicubi esse, et non apparet, quomodo tangant illum locum}).

The post-Tridentine era produced the most exhaustive single treatise on the angels in Catholic history: Francisco Su\'arez's \emph{Tractatus de angelis}, composed before his death in 1617 and published at Lyon in 1620. Drawing freely upon Thomist, Scotist, and patristic arguments, Su\'arez charts an independent synthesis across every contested question. Rejecting both Bonaventure's spiritual matter and Scotus's formally distinct \emph{haecceitas}, Su\'arez teaches that every finite substance is individuated directly by its own intrinsic entity---which leaves God wholly free to create multiple angels in one species, even if Su\'arez grants it probable that each angel in fact differs in degree of perfection. On angelic location he sides closer to Scotus's substantial definitive presence than to Cajetan's purely operative contact; on illumination and speech he modifies Aquinas's exclusion of newly impressed intelligible species; on guardianship he elaborates the roles of personal, national, and ecclesial angels with canonical precision; and on Lucifer's sin (\emph{De angelis}~VII.13) he famously develops the Incarnation-refusal thesis alongside Thomist and Scotist accounts, analyzing how a brilliant spirit could rebel against the revealed decree of the hypostatic union---or even covet that supreme union for himself---without contradicting natural intellectual clarity.

\subsection{The Byzantine Tradition: Cosmic Liturgy and Uncreated Light}

While Western scholasticism honed metaphysical distinctions over individuation, location, and species, the Byzantine Christian East cultivated an unbroken patristic angelology centered on worship, hierarchy, and deifying participation in divine light. That liturgical horizon was already fully formed in the mid-fourth century: a century and a half before Pseudo-Dionysius wrote, St.~Cyril of Jerusalem (\emph{Mystagogical Lecture}~5.6, quoted in full in Section~\ref{sec:patristic-witness}) had already gathered all nine biblical choirs around the altar at the climax of the Eucharistic anaphora, uniting the baptized with the Seraphim's \emph{Sanctus}.

When the Dionysian corpus emerged around the year 500, its vision of the nine celestial choirs as mirrors of divine light descending into the earthly liturgy resonated immediately with this Eastern worship. St.~Maximus the Confessor (c.~580--662), who helped secure the orthodox reception of Pseudo-Dionysius, wove angelic hierarchy directly into the sacramental life of the Church. In his \emph{Mystagogia}, Maximus contemplates the church edifice and the Eucharistic synaxis as a living icon of the whole cosmos and of the human soul: during the Divine Liturgy, the earthly congregation does not merely echo heaven from below, but is lifted by the Holy Spirit into one conciliar choir with the bodiless powers around Christ, joining visible and invisible creation in a single symphony of praise.

In the eighth century, St.~John of Damascus (c.~675--749) gave the Byzantine East its enduring doctrinal standard in Books~II.2--4 of \emph{De fide orthodoxa} (\emph{An Exact Exposition of the Orthodox Faith})---a synthesis equally treasured in Paris, where Aquinas cites it across every question of his treatise. Gathering Cappadocian apophaticism together with ``that most holy, and sacred, and gifted theologian, Dionysius the Areopagite'' (see Section~\ref{sec:patristic-witness}), Damascene locks five pillars into Byzantine angelology: the angels are incorporeal only relative to the transcendent Deity; they act mentally in a single place without bodily extension or bilocation; their incorporeality makes angelic choice irrevocable as death is for men; their nine orders follow the Dionysian three-triad structure in exact descending sequence; and created prior to the visible cosmos (following St.~Gregory of Nazianzus), the chief guardian of the earthly realm fell freely into privation without any evil nature, remaining wholly bounded by God's dispensation.

After the Seventh Ecumenical Council (Nicaea~II, 787) vindicated the veneration of sacred icons of Christ, the Mother of God, the angels, and the saints, Byzantine homiletics and hymnography celebrated the ``bodiless powers'' (\gk{as\={o}matoi dynameis}) with ardent devotion. That homiletic stream is represented in the Greek \emph{Oratio in sanctos angelos} printed as Oration~6 in Migne's \emph{Patrologia Graeca} (vol.~99) under the name of St.~Theodore the Studite (759--826)---an attribution assigned by Cardinal Angelo Mai, who noted that an altered excerpt had also circulated among the spuria of St.~John Chrysostom.

The Byzantine tradition reached its final classical expression in St.~Gregory Palamas (1296--1359). In his \emph{Triads in Defense of the Holy Hesychasts} and his \emph{One Hundred and Fifty Chapters} (\emph{Capita~150}), Palamas addresses the question at the heart of every theology of beatitude: how can finite spirits genuinely see and share in God without erasing the infinite gulf between Creator and creature? Answering through the distinction between God's inaccessible essence (\gk{ousia}) and his uncreated divine energies (\gk{energeiai}), Palamas teaches that not even the Seraphim ever comprehend or merge into the divine essence; yet the holy angels participate directly, eternally, and without intermediary in God's uncreated energies---the very light of glory revealed on Mount Tabor---and transmit that illuminating radiance downward to lower ranks and to deified saints. Approaching from a distinct conceptual grammar than Latin scholasticism's \emph{lumen gloriae} and non-comprehensive vision of the divine essence (\ST{I, q.~12, aa.~5, 7}; \emph{De veritate}~q.~8, a.~2), the Byzantine line guards the same twin truth of Catholic faith: God remains utterly transcendent in his being, yet gives himself truly in deifying light to the angels and saints gathered around his throne.
